%% related_work.tex

\section{Related Work}
\label{sec:related-work}

We position our work along three axes: prior studies of security vulnerabilities in LLM agents, the emerging literature on self-healing and autonomous error recovery in agent systems, and the long tradition of privilege management in operating systems. Our contribution sits at the intersection of these three bodies of work: we show that the \emph{recovery sink}---an artifact of self-healing design---constitutes a privilege-containment gap that prior security work has not examined, and we demonstrate that the classical OS exception-handler privilege invariant ports naturally to close it.

\subsection{Security Vulnerabilities in LLM Agents}
\label{sec:related-injection}

\textbf{Prompt injection.} The canonical vulnerability class for LLM agents is \emph{prompt injection}, in which adversary-controlled text is incorporated into the model's context and subsequently obeyed as if it were a trusted instruction. Greshake et al.\ introduced \emph{indirect prompt injection}, demonstrating that a fetched web page can embed directives that the agent executes with the privileges of its system prompt~\cite{greshake2023indirect}. Willison~\cite{willison2023prompt} catalogued the practical exploitability of injection against production agent frameworks. Liu et al.~\cite{liu2024formalizing} formalized and benchmarked prompt injection attacks and defenses, confirming that injection succeeds with high probability against virtually all production agent architectures.

\textbf{Tool and capability manipulation.} A complementary line of work examines attacks against the \emph{tool layer} and the \emph{capability layer}. Tool-poisoning attacks~\cite{tool-poisoning} demonstrate that malicious MCP servers can manipulate tool descriptions to induce harmful tool calls; tool-slot attacks~\cite{tool-slots} show that adversarial tool outputs can hijack subsequent agent reasoning. CaMeL~\cite{debenedetti2025camel} proposes a capability-based defense that separates control and data flows on the forward path, preventing untrusted data from influencing program flow at tool-call boundaries. These works share a common threat model: the attacker controls content that flows \emph{into} the agent---either as input, as retrieved context, or as tool return values---and the defense is therefore placed at the corresponding entry points (input sanitization~\cite{spotlighting}, instruction-hierarchy enforcement~\cite{instruction-hierarchy}, capability-based data-flow control~\cite{debenedetti2025camel}).

\textbf{Research gap.} Despite the maturity of this literature, every prior attack and defense we are aware of operates on the \emph{forward execution path}: content enters the agent, is processed, and either is obeyed or rejected at the entry boundary. We observe that this framing leaves a structural blind spot. When an agent fails---whether because of injected content or for entirely benign reasons---the recovery pipeline may take a \emph{side-effecting} action (such as replacing the failing agent's role) on the basis of a diagnostic LLM's recommendation, where the diagnostic input was influenced by content that originated outside the trusted computing base. The privilege invariant governing such recovery actions has not been systematically examined. No prior work, to our knowledge, has ported the OS exception-handler privilege rule to the recovery path of an LLM agent system. This paper closes that gap.

Table~\ref{tab:related-comparison} summarizes the specific delta between our contribution and the closest prior work. The key distinctions are: (i)~we target the \emph{recovery sink} (the role-replacement primitive), not the forward execution path; (ii)~our defense is \emph{structural} (a privilege-downgrade invariant enforced at a centralized choke point), not content-based (which is inherently vulnerable to paraphrasing); and (iii)~the defense guarantee is LLM-independent \emph{by construction} (a whitelist check), not probabilistic.

\begin{table}[htbp]
\centering
\small
\caption{Comparison with prior agent-security work. ``Forward path'' = defense at input/tool-entry boundary; ``Recovery sink'' = defense at the role-replacement / side-effect boundary. ``Content-level'' = defense relies on detecting adversarial text; ``Structural'' = defense relies on framework-level invariants (privilege checks) independent of text content.}
\label{tab:related-comparison}
\setlength{\tabcolsep}{3pt}
\begin{tabular}{lcccc}
\toprule
\textbf{Work} & \textbf{Vulnerable Surface} & \textbf{Defense Path} & \textbf{Defense Level} & \textbf{LLM-independent?} \\
\midrule
Greshake et al.~\cite{greshake2023indirect} & Forward (retrieved docs) & Forward & Content-level & No \\
Spotlighting~\cite{spotlighting} & Forward (tool output) & Forward & Content-level & No \\
Instruction Hierarchy~\cite{instruction-hierarchy} & Forward (direct+indirect) & Forward & Content-level & No \\
CaMeL~\cite{debenedetti2025camel} & Forward (tool invocation) & Forward & Structural (capability) & Yes \\
Liu et al.~\cite{liu2024formalizing} & Forward (taxonomy) & Forward & --- & --- \\
\midrule
\textbf{This work} & \textbf{Recovery sink (role replace)} & \textbf{Recovery} & \textbf{Structural (privilege)} & \textbf{Yes (by construction)} \\
\bottomrule
\end{tabular}
\end{table}

\subsection{Self-Healing and Autonomous Error Recovery}
\label{sec:related-selfhealing}

\textbf{Reflection and self-correction.} A substantial body of work improves LLM agent robustness by having the model reflect on its own failures and retry. Reflexion~\cite{reflexion} augments the agent with a verbal reinforcement loop that feeds a self-generated critique into the next iteration; Self-Refine~\cite{self-refine} iterates refinement with self-produced feedback; CRITIC~\cite{critic} externalizes verification to tool-augmented self-critique. These mechanisms uniformly assume that the failure payload (the critique, the error message, the verification trace) is \emph{trustworthy diagnostic signal} to be acted upon, not adversarial content to be defended against.

\textbf{Topology mutation and architectural self-repair.} Multi-agent frameworks additionally restructure the agent graph on failure. AutoGen~\cite{autogen} supports dynamic agent role reassignment; SWE-agent~\cite{swe-agent} reconfigures its tool interface when an action fails; the \system{} system~\cite{omo} and the MetaGPT framework~\cite{meta-agent} describe orchestrators that swap in higher-capability models or replace a failing agent with a different role. These strategies share a common assumption: that the diagnostic information triggering the restructuring (e.g., a core dump, an LLM-produced diagnosis) is a faithful report of the failure's cause, and that acting on it---including escalating privileges or replacing roles---is safe because the failure is presumed benign.

\textbf{Implicit privilege at the recovery sink.} The robustness gains of self-healing are well-documented, but the privilege implications of recovery actions have not been scrutinized. Recovery strategies that restructure the agent graph systematically exercise privileged capabilities: role replacement, model switching, tool reconfiguration. Each of these is a privileged operation that, on the forward execution path, would be gated by a permission checker. On the recovery path, the same operations are executed on the basis of a diagnostic LLM's recommendation, whose input (the core dump) may embed attacker-influenced content. Our work is the first to formalize this as a privilege-containment failure at the recovery sink---an instance of CWE-862 and the Confused Deputy problem---and to propose a structural fix.

\textbf{Concurrent and adjacent work (novelty scope).} We conducted a systematic search for prior work targeting privilege enforcement on the recovery path in LLM agents, using the keywords ``recovery path privilege'', ``role replacement authorization'', ``agent recovery confused deputy'', ``topology mutation security'', and ``exception handler privilege LLM'' across arXiv (2023--2026), Google Scholar, and the major software-engineering and security venues. We found no prior work that (i)~formalizes a privilege-downgrade invariant for LLM agent recovery actions, (ii)~connects the recovery sink to the Confused Deputy problem, or (iii)~proposes a centralized reauthorization gate for side-effecting recovery strategies. The closest adjacent works we are aware of operate on the \emph{forward} path and are complementary to ours. Zhan et al.~\cite{zhan2025adaptive} evaluate 8 defenses against indirect prompt injection and bypass all of them with adaptive attacks; their finding that content-level forward-path defenses fall to adaptive attacks reinforces our decision to make the V2 defense structural (a programmatic whitelist check), not content-based. CaMeL~\cite{debenedetti2025camel} proposes a capability-based defense that separates control and data flows on the forward path; our \reauthgate{} ports the same capability-based principle to the \emph{recovery} path, suggesting the two defenses are complementary (CaMeL guards the forward execution; \reauthgate{} guards the recovery side effects). The OWASP Top 10 for LLM Applications~\cite{owasp-llm-top10} enumerates \emph{LLM01: Prompt Injection} and \emph{LLM06: Excessive Agency} as forward-path risks but does not classify the recovery sink as a distinct privilege boundary. We therefore believe our contribution is novel within the LLM-agent security literature, while acknowledging that the field is moving rapidly and concurrent work may emerge between submission and publication.

\subsection{Privilege Management in OS and AI}
\label{sec:related-privilege}

\textbf{OS privilege invariants.} Classical operating systems enforce a small set of invariants that have proven robust over decades. The \emph{principle of least privilege}~\cite{saltzer1975protection} holds that every component should operate with the minimum privileges required for its task; the \emph{exception-handler privilege invariant} holds that an exception handler executes with the privileges of the faulting context, never higher (e.g., a userspace page fault is handled in kernel mode but cannot escalate the faulting process's privileges); and \emph{capability-based authority}~\cite{capabilities} holds that a privileged deputy must present a capability that justifies the action, not merely act on behalf of a request. Sandboxing~\cite{capsules,seccomp} extends these invariants by restricting the ambient authority of a process to a declared subset.

\textbf{Privilege in AI systems.} The AI security literature has begun importing these principles. \cite{instruction-hierarchy} enforces an instruction-hierarchy that treats system prompts as higher-privilege than user content, an analog of kernel/userspace separation; CaMeL~\cite{debenedetti2025camel} extracts control and data flows from a trusted query and uses capabilities to prevent untrusted data from influencing program flow at tool-call boundaries, an analog of capability gating. These works port the \emph{forward-path} versions of the OS invariants (input validation, capability gating) to the agent setting.

\textbf{The Confused Deputy problem.} Our attack vector is a textbook instance of the \emph{confused deputy} problem, originally identified by Hardy~\cite{hardy1988confused}: a privileged program (in our setting, the recovery orchestrator operating with role-replacement authority) is tricked by a less-privileged caller (the attacker-controlled external content embedded in the core dump) into misusing its authority on the caller's behalf. The recovery pipeline holds the privilege to swap agent roles---a capability not granted to ordinary tool calls---and, in the baseline configuration, it exercises this capability on the basis of an LLM diagnosis whose input was supplied by attacker-influenced content. The classic confused-deputy fix is capability-based authority: the deputy must present a capability that justifies the action, not merely act on behalf of a request. Our \reauthgate{} is precisely such a capability check: it demands that the proposed role replacement be in the pre-approved downgrade set of the failing role, regardless of what the diagnostic LLM suggested. The confused-deputy framing also clarifies why the fix must be \emph{structural} rather than content-based: the deputy's vulnerability lies in \emph{whose authority} it is exercising, not in the textual content of the request, so no content-level defense (sanitization, framing, instruction hierarchy) can close it. To our knowledge, this is the first explicit connection between the confused-deputy problem and LLM agent recovery channels, and it positions our work within a long-standing line of capability-based security work~\cite{capabilities,hardy1988confused}.

\textbf{Cross-cutting concern literature.} The software-engineering literature on cross-cutting concerns~\cite{aspectj} motivates our central design choice: permission enforcement on the recovery path is an aspect that every side-effecting strategy needs, but is orthogonal to the strategy's core logic (failure diagnosis and recovery planning). Aspect-oriented programming and middleware-based enforcement in classical distributed systems both centralize such cross-cutting concerns at well-defined boundaries. Our \reauthgate{} applies the same principle to the recovery orchestrator: the security property ``no recovery action exceeds the privilege of the failed execution'' becomes a local, inspectable property of one orchestrator method, rather than an emergent property of fourteen independently-authored strategies.

\textbf{Our contribution: porting the reverse-path privilege invariant.} The observation that motivates this paper is that the OS exception-handler privilege invariant applies not only on the forward execution path but also---and critically---on the \emph{reverse} recovery path. Classical kernels learned long ago that exception handlers must not escalate privileges; LLM agent systems have not yet learned the corresponding lesson. Our defense is a direct port of this invariant:

\begin{itemize}[leftmargin=*]
    \item The \reauthgate{} is the agent analog of the exception-handler privilege invariant. It enforces $P(\text{recovery action}) \leq P(\text{failed execution})$, so a recovery strategy may remediate a failure but never escalate beyond it---just as a kernel exception handler services a fault without granting the faulting process new privileges. The same principle underlies the confused-deputy fix~\cite{hardy1988confused}: the privileged deputy may not exercise authority not granted to the caller, even when the caller's request looks superficially legitimate.
\end{itemize}

The mechanism is lightweight (single record allocation, $O(1)$ role-set lookup; see Section~\ref{sec:evaluation}), framework-agnostic in principle (it depends only on the existence of a recovery orchestrator and a permission checker, not on any specific agent architecture), and---critically---\emph{local}: it is implemented at a single well-defined choke point rather than scattered across the recovery strategy layer. This locality is itself an OS principle: classical kernels enforce privilege invariants at well-defined boundaries (syscall entry, exception dispatch), not inside each driver or subsystem.

In this sense, our work is not a novel cryptographic mechanism but a \emph{transfer} of a proven OS invariant to a domain where it has been conspicuously absent. The validation result---that the transfer closes the privilege-containment gap with zero false alarms and sub-millisecond machinery overhead---suggests that the gap was one of \emph{attention}, not of mechanism: the recovery sink was overlooked, not undefendable.
